\section{Introduction}

The exploration of deep space during upcoming lunar and Martian expeditions demands unprecedented vigilance regarding spacecraft environmental microbiology and life support systems integrity \cite{nickerson2000microgravity}. Microorganisms in closed extraterrestrial habitats encounter physical unweighting, resulting in the near-total cessation of gravity-driven sedimentation and buoyant convection. Under these quiescent, low-shear fluid regimes, mass transport around single-celled organisms becomes dominated strictly by molecular diffusion, altering boundary layer mechanics, nutrient gradients, and mechanical membrane tension \cite{nickerson2000microgravity,falkinham2015common}. Over decades of flight investigations aboard the Space Shuttle and the International Space Station (ISS), opportunistic pathogens including \textit{Pseudomonas aeruginosa} (OSD-14/15/554), \textit{Salmonella enterica} (OSD-11/277), and \textit{Staphylococcus aureus} (OSD-145/683) have consistently exhibited enhanced virulence phenotypes, accelerated biofilm architecture formation, and elevated antimicrobial tolerance.

Among persistent biofilm formers isolated from terrestrial water distribution systems and spaceflight hardware, nontuberculous mycobacteria (NTM) represent a profound concern \cite{falkinham2015common}. \textit{Mycobacterium marinum} is a natural pathogen of ectotherms and human cutis, sharing substantial genomic synteny, specialized secretion systems (Type VII secretion), and cell wall lipid complexity with obligate human pathogens such as \textit{Mycobacterium tuberculosis} and opportunistic pathogens like \textit{Mycobacterium avium}. Because \textit{M. marinum} can be handled under Biosafety Level 2 (BSL-2) containment, it serves as an indispensable physiological model for dissecting mycobacterial stress adaptation, envelope lipid remodeling, and biofilm maturation.

In a landmark ground-based simulation study deposited in the NASA Open Science Data Repository (OSDR) under accession \textbf{OSD-528} (GLDS-528), Clary, Harrison, and colleagues \cite{clary2022development} characterized the global response of \textit{M. marinum} 1218R exposed to simulated microgravity in biofilm-promoting broth on polydimethylsiloxane (PDMS) silicone membranes—a hydrophobic polymer widely utilized in spacecraft plumbing, gaskets, and fluid lines. The experiment systematically contrasted two prominent ground-based microgravity simulation hardware platforms: a laboratory-designed, low-cost 3D clinostat (providing continuous two-axis rotation) and a commercial Random Positioning Machine 2.0 (RPM 2.0, providing multi-axis random velocity vectoring), against static 1g ground controls housed in the same incubator shelf ($31^\circ$C, 4 days). Clary et al. demonstrated that both simulators produced similar biofilm structures, comparable reductions in planktonic suspension growth, and concordant increases in rifampicin tolerance \cite{clary2022development}.

Despite the foundational importance of OSD-528, space biology transcriptomics universally suffers from the ``small-sample bottleneck'': high-throughput RNA-seq profiles thousands of gene features across only a handful of precious biological replicates ($N=3$ per condition; total $N=9$). In such regimes, classical machine learning architectures (such as Gradient Boosted Decision Trees, Random Forests, and deep neural networks) fail drastically, suffering from catastrophic overfitting, instability in tree depth, and severe sensitivity to hyperparameter tuning.

Recently, Hollmann et al. (\textit{Nature} 2025) introduced **TabPFN** \cite{hollmann2025accurate}, a tabular foundation model trained via in-context learning (ICL) on millions of synthetic causal structural priors. TabPFN delivers instant, single-forward-pass inference tailored specifically for small- to medium-sized datasets ($N \le 10,000$ samples, $P \le 500$ features), completely eliminating hyperparameter tuning while dominating tree-based ensembles across small tabular benchmarks.

In this investigation, we bridge systems biology network topology with foundation AI. We deploy Weighted Gene Co-expression Network Analysis (WGCNA) \cite{langfelder2008wgcna} to reduce the 5,424-gene space of \textit{M. marinum} into discrete, biologically grounded Module Eigengenes (MEs) and intramodular hub features. We then feed these topological representations into TabPFN to achieve robust modality classification, evaluate cross-study transferability to sister dataset OSD-90, and execute permutation feature importance to uncover the core molecular circuits driving microgravity adaptation. Finally, we package the complete research object in accordance with FAIR data principles \cite{wilkinson2016fair} with RO-Crate v1.1 and Zenodo schemas, delivering an end-to-end reproducible pipeline.
